\documentclass[11pt]{article}
\usepackage[utf8]{inputenc}
\usepackage[T1]{fontenc}
\usepackage{amsmath,amssymb}
\usepackage{graphicx}
\usepackage{booktabs}
\usepackage{hyperref}
\usepackage[margin=1in]{geometry}
\usepackage{natbib}
\bibliographystyle{unsrtnat}

\title{Disseminating cells in human oral tumours acquire an EMT cancer stem cell state that is predictive of metastasis}

\author{
Adrian Biddle$^{1,*}$ \and Luke Gammon$^{1}$ \and Ian C.\ Mackenzie$^{1}$\\
\\
\small $^{1}$Centre for Cell Biology and Cutaneous Research, Blizard Institute,\\
\small Barts and The London School of Medicine and Dentistry, Queen Mary University of London, UK\\
\small $^{*}$Corresponding author
}

\date{}

\begin{document}
\maketitle

\begin{abstract}
Cancer stem cells (CSCs) undergo epithelial-mesenchymal transition (EMT) to drive metastatic dissemination in experimental cancer models. However, tumour cells undergoing EMT have not been observed disseminating into the tissue surrounding human tumour specimens, leaving the relevance to human cancer uncertain. We have previously identified both EpCAM and CD24 as markers of EMT CSCs with enhanced plasticity. This afforded the opportunity to investigate whether retention of EpCAM and CD24 alongside upregulation of the EMT marker Vimentin can identify disseminating EMT CSCs in human tumours. Examining disseminating tumour cells in over 12,000 imaging fields from 84 human oral cancer specimens, we see a significant enrichment of single EpCAM, CD24 and Vimentin co-stained cells disseminating beyond the tumour body in metastatic specimens. Through training an artificial neural network, these predict metastasis with high accuracy (cross-validated accuracy of 87--89\%). In this study, we have observed single disseminating EMT CSCs in human oral cancer specimens, and these are highly predictive of metastatic disease.
\end{abstract}

\section{Introduction}

In multiple types of carcinoma, cancer stem cells (CSCs) undergo epithelial-mesenchymal transition (EMT) to enable metastatic dissemination from the primary tumour \citep{biddle2011,lambert2017}. This model of metastatic dissemination has been built from studies using murine models and human cancer cell line models. However, this process has not been observed in human tumours in the in vivo setting, leading to uncertainty over the relevance of these findings to human tumour metastasis \citep{bill2015,williams2019}.

A key complication with efforts to study metastatic processes in human tumours is the inability to trace cell lineage. As cancer cells exiting the tumour downregulate epithelial markers whilst undergoing EMT, they become indistinguishable from the mesenchymal non-tumour cells surrounding the tumour \citep{li2016}. Cytokeratins, which are commonly used to identify epithelial tumour cells, are lost during EMT and therefore cannot mark disseminating cells that have undergone this transition.

We have previously shown that CSCs in oral squamous cell carcinoma (OSCC) switch between epithelial and EMT states, with EMT CSCs exhibiting enhanced phenotypic plasticity and therapeutic resistance \citep{biddle2011,biddle2016}. Critically, we identified that while epithelial keratins are lost during EMT, the surface marker EpCAM is retained \citep{biddle2016}. Combined with CD24 as a marker of plastic EMT CSCs and Vimentin as a mesenchymal marker, this creates a three-marker combination---EpCAM$^+$Vim$^+$CD24$^+$---that can identify tumour cells that have undergone EMT.

Here, we examine whether this marker combination can detect disseminating EMT CSCs in human OSCC specimens and test whether the presence of these cells predicts metastatic status.

\section{Results}

\subsection{EpCAM, Vimentin and CD24 co-staining identifies EMT CSCs}

We first validated that the EpCAM$^+$Vim$^+$CD24$^+$ staining profile identifies EMT CSCs using the CA1 cell line and its EMT-stem sub-line. EpCAM$^+$Vim$^+$CD24$^+$ cells were greatly enriched in the EMT-stem sub-line, comprising 41\% of the population compared to 2.1\% in the CA1 parental line (Figure~1A, B, E). Cells with this profile were absent from normal keratinocyte and cancer-associated fibroblast cultures. In contrast, pan-keratin$^+$Vim$^+$CD24$^+$ staining showed very little signal and no enrichment in the EMT-stem sub-line (Figure~1C, D, E), confirming that epithelial keratins are lost but EpCAM is retained during EMT.

Imaging human OSCC specimens, we detected single cells co-expressing EpCAM, Vimentin and CD24 in the stromal region surrounding the tumour (Figure~1F), confirming that these cells can be detected in human tumour tissue.

\subsection{EpCAM$^+$Vim$^+$CD24$^+$ cells are enriched in the stroma of metastatic tumours}

We stratified 24 human primary OSCC specimens into 12 tumours with evidence of lymph node metastasis or perineural spread and 12 that remained metastasis free. Single cells co-expressing EpCAM, Vimentin and CD24 were abundant in the stroma surrounding metastatic tumours but not in non-metastatic tumours or normal epithelial regions (Figure~2A--C). In contrast to EpCAM, pan-keratin staining did not identify cells in the stroma surrounding metastatic tumours (Figure~2D).

Using an automated 4-colour immunofluorescent imaging and analysis protocol, we quantified expression of EpCAM, Vimentin and CD24 for every nucleated cell across 3,500 manually curated imaging fields from 24 tumours containing 2,640,000 total nucleated cells. EpCAM$^+$Vim$^+$CD24$^+$ cells were enriched in the stroma compared to the tumour body, with significantly greater accumulation in metastatic tumour stroma compared to non-metastatic tumour stroma (Figure~2F). Notably, EpCAM$^+$Vim$^+$CD24$^-$ cells were also enriched in the stroma but showed no difference between metastatic and non-metastatic tumours, highlighting the requirement for the plasticity marker CD24 in identifying disseminating metastatic CSCs.

A further 60 tumours, evenly stratified on the same criteria, displayed the same pattern: individual disseminating EpCAM$^+$Vim$^+$CD24$^+$ cells were observed in metastatic tumours only, across over 9,000 imaging fields at the tumour-stroma boundary.

\subsection{Confirmation of tumour cell identity}

To exclude the possibility that observed EpCAM$^+$Vim$^+$CD24$^+$ cells in the stroma are non-tumour cell types, we analysed a published single-cell RNA-seq dataset for human head and neck cancer \citep{puram2017}. Of tumour cells, 12\% (267/2215) were EpCAM$^+$Vim$^+$CD24$^+$, whereas only 0.8\% (29/3687) of non-tumour cells showed this profile. EpCAM is a specific epithelial marker not expressed in stromal or immune cells \citep{keller2019}, further supporting that the observed cells are a tumour cell population.

\subsection{Machine learning prediction of metastasis}

Starting with the EpCAM, Vimentin and CD24 immunofluorescence grey levels for each nucleated cell, we used a supervised machine learning approach to predict whether an imaging field came from a metastatic or non-metastatic tumour (Figure~3A). Using 3,500 imaging fields from the first batch of 24 tumours, we compared several machine learning classifiers \citep{pedregosa2011}.

The best performing classifiers for EpCAM, Vimentin and CD24 were the artificial neural network (ANN) \citep{rumelhart1986} and support vector machine (SVM) \citep{smola2004}, with 10-fold cross-validated F1 accuracy scores of 91\% and 87\%, respectively (Figure~3B). For the ANN, the area under the curve (AUC) accuracy score was 87\%, with 94\% sensitivity and 82\% specificity. Training with pan-keratin, Vimentin and CD24 gave substantially worse prediction across all classifiers (Figure~3C), consistent with the absence of pan-keratin$^+$ cells disseminating in the stroma.

To validate this finding, over 9,000 imaging fields at the tumour-stroma boundary from the second batch of 60 tumour specimens (containing over 8.5 million nucleated cells) were fed into an ANN classifier trained and tested independently of the first batch. This showed good alignment between training and validation sets and achieved an 89\% accuracy score after 12 training epochs (Figure~3D).

\section{Discussion}

In this study, we have demonstrated the presence of single disseminating EMT CSCs in human oral cancer specimens. These cells are identified by co-expression of EpCAM, Vimentin and CD24, a staining profile that marks EMT CSCs with enhanced plasticity. Their specific enrichment in the stroma surrounding metastatic tumours, but not non-metastatic tumours, provides direct human tumour evidence supporting the experimental model of EMT-driven metastatic dissemination \citep{biddle2011,tsai2012,ocana2012}.

The retention of EpCAM during EMT is critical to this approach. While conventional epithelial markers such as cytokeratins are lost as cells undergo EMT, making disseminating tumour cells indistinguishable from surrounding stromal cells, EpCAM is maintained \citep{biddle2016,keller2019}. This provides a tumour lineage marker that persists through the mesenchymal transition. The additional requirement for CD24 co-expression to distinguish metastatic from non-metastatic dissemination highlights the importance of the CSC plasticity phenotype in the metastatic process \citep{biddle2011,liu2014,lawson2015}.

The ability to predict metastatic status with 87--89\% accuracy using an ANN classifier trained on three-marker immunofluorescence data has potential clinical utility. OSCC is one of the top ten cancers worldwide, and lymph node metastasis is the single most important predictor of outcome \citep{sano2007}. Previous machine learning approaches using cytokeratin immunohistochemistry, clinicopathological data and serum biomarkers have achieved AUCs of 75--82\% in various cancer types \citep{tseng2019,bur2019,takamatsu2019}. Our approach achieves comparable or superior accuracy using a mechanistically motivated marker combination.

\section{Materials and Methods}

\subsection{Cell culture}
The CA1 OSCC cell line and its EMT-stem sub-line were cultured as previously described \citep{biddle2016}. Normal keratinocytes and cancer-associated fibroblasts were cultured as controls.

\subsection{Immunofluorescent staining}
Tumour specimens were stained with four colours: EpCAM (yellow), Vimentin (red), CD24 (green), and DAPI nuclear stain (blue). An alternative staining panel substituted pan-keratin for EpCAM. Automated whole-slide imaging was performed and grey level intensities for each marker were obtained for every nucleated cell in each imaging field.

\subsection{Tumour specimen selection}
A total of 84 human OSCC specimens were analysed in two batches: 24 specimens (12 metastatic, 12 non-metastatic) and 60 specimens (30 metastatic, 30 non-metastatic). Metastatic status was defined by evidence of lymph node metastasis or perineural spread.

\subsection{Image segmentation}
Imaging fields were segmented to separate the tumour body from adjacent stroma. In the first batch, 3,500 imaging fields were manually curated. In the second batch, over 9,000 imaging fields at the tumour-stroma boundary were analysed using a variation of the segmentation protocol.

\subsection{Machine learning}
Grey level intensities for each marker in each nucleated cell were used as features for supervised classification. Several algorithms were compared: artificial neural network (ANN), support vector machine (SVM), naive Bayes \citep{rish2005}, random forest \citep{deparis2009}, and k-nearest neighbour \citep{benhamou1975}. Performance was assessed using 10-fold cross-validated F1 scores. The scikit-learn library \citep{pedregosa2011} was used for implementation.

\subsection{Single-cell RNA-seq analysis}
A published scRNAseq dataset for human head and neck cancer \citep{puram2017} was analysed for co-expression of EpCAM, Vimentin and CD24 in tumour versus non-tumour cell populations.

\bibliography{references}

\end{document}
